\documentclass[10pt]{article}
\usepackage[a4paper,margin=0.5in]{geometry}
\usepackage{amsmath,amssymb,amsfonts,mathtools}
\usepackage[dvipsnames]{xcolor}
\usepackage{array,longtable,pifont,booktabs,float}
\usepackage[colorlinks=true,linkcolor=blue,citecolor=blue]{hyperref}
\allowdisplaybreaks
\numberwithin{equation}{section}
\setlength{\jot}{2pt}
\setlength{\LTpre}{4pt}\setlength{\LTpost}{4pt}
\setlength{\emergencystretch}{3em}
\newcommand{\KK}{\mathcal K}
\newcommand{\NN}{\mathcal N}
\newcommand{\lab}[1]{\texttt{#1}}
\newcommand{\ok}{\textcolor{OliveGreen}{\ensuremath{\checkmark}}}
\newcommand{\bad}{\textcolor{red}{\ding{55}}}
\def\mf#1{\mathbb{#1}}
\setcounter{tocdepth}{2}

\title{The three-$\rho$ terms proportional to $\log(\vee/\Lambda)$:\\ check of the list, fixes, and the two-$\rho$ terms}
\date{}

\begin{document}
\maketitle
\tableofcontents

%==========================================================================================
\section{Question and answer}\label{sec:answer}

\paragraph{Question.} You sent all three-$\rho$ terms proportional to $\log(\vee/\Lambda)$, with their c.c.:
\begin{itemize}
\item the two remainders $d^3N_a$, $d^3N_b$ of the four-$\rho$ terms;
\item the genuine three-$\rho$ terms of Group I Rows II--V, Group II Row I and Group III Rows I, II, III, V, VI.
\end{itemize}
In total there are 160 terms, written in common variables with kernel classes $\alpha,\alpha',\beta,\beta',\gamma,\gamma'$. You asked three things:
\begin{enumerate}
\item Do they add up to zero after reordering (symmetrizing)?
\item If not, fix them so that they do.
\item Give the corresponding two-$\rho$ terms, which go to JIMWLK and BFKL.
\end{enumerate}

\paragraph{Answer.}
\begin{enumerate}
\item \textbf{As sent, the terms do not cancel} (Sec.~\ref{sec:sent}).
\begin{itemize}
\item Pointwise in $z$, all 36 kernel monomials are left over.
\item Grouped into colour networks with exact coefficients, 24 of 50 networks do not cancel.
\item The list agrees with the earlier transcription, as written, except for three typos. So the two fixes found before are still needed.
\end{itemize}
\item \textbf{Five changes make them cancel exactly} (Sec.~\ref{sec:fixes}). After them, all 40 colour networks cancel with exact coefficients, and the pointwise residual is exactly zero (Table~\ref{tab:states}).
\begin{itemize}
\item[T1] Group III, Row I, term 15 ($\beta'$): $\rho^h(x')\to\rho^c(x')$. As written, the index $h$ occurs three times. Term 11 of the same row (its $\alpha'$ partner) has $\rho^c(x')$.
\item[T2] Group III, Row II, term 13 ($\beta'$): $U^{hl}(y)\to U^{hl}(x)$. Term 9 (its $\alpha'$ partner) has $U^{hl}(x)$.
\item[T3] Group I, Row II, part 1, the eight $\alpha$ terms: the kernel is written $\KK(x-w)\cdot\KK(x'-w')\,\KK(x'-w)\cdot\KK(y-z)$. The $\alpha$ kernel is $\KK(x-w)\cdot\KK(x'-w')\,\KK(x-z)\cdot\KK(y-z)$: $\KK(x'-w)\to\KK(x-z)$.
\item[F4] Group I, Row II: remove the $\alpha$ terms of the $U(w)\mathbb A^{(3)}$ sub-terms. These are part 1, terms 1, 2, 5, 6, and their c.c., part 2, terms 1, 2, 5, 6. This part, $\mathfrak F_\alpha$, is the commutator $[\mathbb N^{(2)},\mathbb A^{(1)}]$. With the $p^+$ ordering of the virtual factor it multiplies $-\log(\vee/k^+)$, not $\log(\vee/\Lambda)$.
\item[F3] Group I, Row III: remove the $\gamma,\gamma'$ terms (part 2, edge kernel $\phi$, terms 9--16) and add $\mathcal D$. $\mathcal D$ has eight terms of class $\beta,\beta'$, written in your notation in Sec.~\ref{sec:F3}.
\end{itemize}
T1--T3 are typos. F3 and F4 are the two fixes of the earlier three-$\rho$ note.
\item \textbf{The two-$\rho$ terms need the operator order of the charges, and your list does not have it} (Sec.~\ref{sec:order}). Your list gives each term in the charge order of the symmetrized check. That is enough for the three-$\rho$ check, but not for the two-$\rho$ terms. There are two differences:
\begin{itemize}
\item[(a)] \emph{c.c.\ terms} (the primed classes): the c.c.\ of an operator product reverses the order of the charges. $ABC$ and $CBA$ have the same symmetrized part but two-$\rho$ terms of opposite sign. Your c.c.\ terms keep the order of the original term.
\item[(b)] \emph{Rows with an anticommutator} in the notes (Group I Row V, Group II Row I, Group III Row III): your list writes $\{B,C\}A$ as $2\,BCA$. The two have the same symmetrized part, but their two-$\rho$ terms differ by $\tfrac12\{[B,C],A\}$. In Group I Row II the anticommutator is antisymmetrized by $f^{abc}$, and your two orderings are exactly $\{\rho,\rho\}$.
\end{itemize}
Reordering your list literally gives wrong two-$\rho$ terms in 11 of the 12 sources (Table~\ref{tab:literal}). Sec.~\ref{sec:oplist} gives the corrected list again, term by term, with your colour indices, in operator order; the changed terms are blue. Its two-$\rho$ terms equal those of the original expressions exactly, row by row.
\item \textbf{The two-$\rho$ terms} are listed source by source in Sec.~\ref{sec:two}. Added to the genuine two-$\rho$ terms of \texttt{Evolution\_two\_rho.tex} and to the lists of \texttt{4rho.tex}, they give JIMWLK$\otimes$LO exactly: @@NS@@ kernel structures, all @@NB@@ independent colour structures with a nonzero entry, exact coefficients (Sec.~\ref{sec:sum}).
\item \textbf{BFKL.} BFKL$\otimes$LO is the coefficient of $\log\vee$. To get it, the $\log(\vee/k^+)$ terms of each row are added to these $\log(\vee/\Lambda)$ terms; the result is BFKL$\otimes$LO exactly. That part is unchanged by this check and is written out in \texttt{TwoRho\_all.pdf}, Sec.~9.
\end{enumerate}

%==========================================================================================
\section{Conventions}\label{sec:conv}
\paragraph{Units.} $\NN\equiv\frac{1}{(2\pi)^3}\frac{g^4}{16\pi^5}\frac1{k^+}\log\frac\vee\Lambda$, so
\begin{itemize}
\item your prefactor $\frac{1}{(2\pi)^3}\frac{ig^4}{16\pi^5}\frac1{k^+}\log\frac\vee\Lambda$ is $i\NN$;
\item $\mp\frac{i}{(2\pi)^3}\frac{g^4}{8\pi^3}\frac{\log(\vee/\Lambda)}{2\pi^2}\frac1{k^+}$ is $\mp i\NN$.
\end{itemize}
Every block is $(\text{prefactor})\int e^{-ik\cdot(w'-w)}[\dots]$ with all positions except $w,w'$ integrated.

\paragraph{Kernel classes} (your labels). With $\KK(X)=X/X^2$,
\begin{align*}
K_\alpha&=\KK(x'-w')\cdot\KK(x-w)\;\KK(x-z)\cdot\KK(y-z), & K_{\alpha'}&=\KK(x'-w')\cdot\KK(x-w)\;\KK(x'-z)\cdot\KK(y-z),\\
K_\beta&=\KK(x-w)\cdot\KK(x'-w')\;\KK(w-z)\cdot\KK(y-z), & K_{\beta'}&=\KK(x-w)\cdot\KK(x'-w')\;\KK(w'-z)\cdot\KK(y-z),\\
K_\gamma&=\KK(x'-w')\cdot\KK(y-z)\;\phi_w(x), & K_{\gamma'}&=\KK(x-w)\cdot\KK(y-z)\;\phi_{w'}(x'),
\end{align*}
with $\phi_w(x)$ the edge kernel of Group I, Row III, part 2. The charge tied to $w$ is $x$, the one tied to $w'$ is $x'$, and the third charge is $y$. In the lists below, term numbers are those of your list, block by block.

\paragraph{Symmetrization and two-$\rho$ terms.} For three charges $A,B,C$, with $S(ABC)$ the fully symmetrized product,
\begin{align}
ABC&=S(ABC)+\tfrac14\big(\{[A,B],C\}+\{[A,C],B\}+\{[B,C],A\}\big)+(\text{one }\rho),\label{eq:plain}\\
A\{B,C\}&=2S(ABC)+\tfrac12\{[A,B],C\}+\tfrac12\{[A,C],B\}+(\text{one }\rho),\label{eq:Aanti}\\
\{B,C\}A&=2S(ABC)-\tfrac12\{[A,B],C\}-\tfrac12\{[A,C],B\}+(\text{one }\rho),\label{eq:antiA}
\end{align}
with $[\rho^a(x),\rho^b(y)]=if^{abc}\rho^c(x)\delta^{(2)}(x-y)$. The three-$\rho$ check uses only $S$; the two-$\rho$ terms are the anticommutators. Everything uses $U^{bc}=U^{cb}$.

%==========================================================================================
\section{The list as sent}\label{sec:sent}
Each term is symmetrized (charges commute). The sum is evaluated pointwise at random positions, random symmetric SU(3) Wilson lines and random charges. It is sorted by kernel monomial, with the charges placed on three slots in all $3!$ ways. Table~\ref{tab:states} shows the result as sent and after each fix.
\begin{table}[H]
\centering\small
\begin{tabular}{@{}l r r l c@{}}
state & terms & kernel monomials & non-cancelling (3 configurations) & largest residual\\ \hline
@@STATES@@
\end{tabular}
\caption{Pointwise test of the symmetrized three-$\rho$ terms. ``As sent'' already reads the $\alpha$ kernel of Group I Row II part 1 as $K_\alpha$ (T3); with the kernel as written, those eight terms match nothing. Only the last line cancels.}
\label{tab:states}
\end{table}
\paragraph{Row by row.} Each row of the list was compared with the earlier transcription of \texttt{Evolution\_three\_rho.tex} and of the four-$\rho$ remainders (the ThreeRho\_cancellation note), term by term.
\begin{itemize}
\item All rows agree exactly with the earlier list ``as written'', except Group III Rows I and II, which differ by the single terms T1 and T2.
\item Row II contains the full $\mathfrak F$ in all four sub-terms, and Row III contains part 2 ($\gamma,\gamma'$). So neither F3 nor F4 is in the list.
\end{itemize}
The earlier note found that these two fixes are needed and sufficient. The colour-network tables of Sec.~\ref{sec:networks} show exactly which networks they repair.

%==========================================================================================
\section{The fixes}\label{sec:fixes}
\subsection{T1, T2: two index typos in Group III}
@@TYPOS@@
In T1 the index $h$ appears three times as written (in $U^{hl}(w)$, $U^{ha}(w')$, $\rho^h(x')$). The free index of the charge at $x'$ must be $c$, the third index of $f^{abc}$, as in the $\alpha'$ term 11 of the same row. In T2 the Wilson line $U^{hl}$ belongs to the charge $\rho^l(x)$, as in the $\alpha'$ term 9.

\subsection{T3: the \texorpdfstring{$\alpha$}{alpha} kernel of Group I, Row II, part 1}
All eight $\alpha$ terms of part 1 are written with $\KK(x-w)\cdot\KK(x'-w')\,\KK(x'-w)\cdot\KK(y-z)$. The Fourier transform of Row II gives $\mathfrak F^i=\KK^i(x-w)\int_z[K(w,y,z)-K(x,y,z)]-(x\leftrightarrow y)$. Its $\alpha$ part is $K(x,y,z)=\KK(x-z)\cdot\KK(y-z)$, and the label $\alpha$ says the same. So the kernel is $K_\alpha$; the $\alpha'$ terms of part 2 (the c.c.) have the correct $K_{\alpha'}$.

\subsection{F4: Group I, Row II}
Split $\mathfrak F^i=\mathfrak F^i_\beta+\mathfrak F^i_\alpha$, with $\mathfrak F^i_\alpha=-\big[\KK^i(x-w)-\KK^i(y-w)\big]\int_zK(x,y,z)$.
\begin{itemize}
\item In the $U(w)\mathbb A^{(3)}$ sub-terms ($U^{de}(x')U^{ad}(w)$ and $U^{de}(w')U^{ad}(w)$), $\mathfrak F_\alpha$ is the commutator $[\mathbb N^{(2)},\mathbb A^{(1)}]$ of the normalization with the LO amplitude (combined note, Part I).
\item Ordering the virtual factor in $p^+$ puts its soft part to the left of $\mathbb A^{(1)}$ and its hard part to the right. Then this commutator comes with $-\log(\vee/k^+)$, not with $\log(k^+/\Lambda)$, so it is not a $\log(\vee/\Lambda)$ term.
\item The $\mathfrak F_\alpha$ part of the $\bar{\mathbb A}^{(3)}$ sub-terms stays.
\end{itemize}
In your list this removes:
@@F4LINES@@

\subsection{F3: Group I, Row III}\label{sec:F3}
Part 2 of Row III (the edge kernel $\phi$, classes $\gamma,\gamma'$) has nothing to cancel against. It is replaced by
\begin{align}
\mathcal D={}&-2i\NN\int e^{-ik\cdot(w'-w)}\,\KK(x'-w')\cdot\KK(y-w)\,\KK(x-z)\cdot\KK(w-z)\notag\\
&\times f^{abd}\big[U^{ab'}(x')-U^{ab'}(w')\big]\rho^{b'}(x')\,U^{de}(x)\rho^e(x)\big[U^{bc}(w)-U^{bc}(y)\big]\rho^c(y)+{\rm c.c.}
\end{align}
\begin{itemize}
\item $\mathcal D$ is Row III in the region where the gluon of $\mathbb A^{(1)}(p^+,z)$ carries almost all of $k^+$ and becomes the measured gluon. There $U(z)\to U(w)$.
\item The charge $y$ emits the measured gluon at LO ($\KK(y-w)$), and a soft gluon at $z$ is exchanged between the measured gluon and $U(x)\rho(x)$ ($\KK(x-z)\cdot\KK(w-z)$).
\item In common variables ($x\leftrightarrow y$) it is of class $\beta$, and its c.c.\ of class $\beta'$.
\end{itemize}
In your list this removes:
@@F3LINES@@
and adds, in your notation, the eight terms of $\mathcal D$:
@@DLINES@@
For the two-$\rho$ terms of $\mathcal D$ the operator order is that of Row III, $\rho(x')\rho(x)\rho(y)$ before the renaming $x\leftrightarrow y$. In common variables it is $\rho(x')\rho(y)\rho(x)$, and the c.c.\ is reversed.

\subsection{The corrected list}\label{sec:corrected}
The full list after the fixes, in your notation and in your (symmetrized) charge order.
\begin{itemize}
\item Removed terms are red: \textcolor{red}{[F3]}, \textcolor{red}{[F4]}.
\item Corrected and added terms are blue: \textcolor{blue}{[T1]}, \textcolor{blue}{[T2]}, \textcolor{blue}{[T3]}, \textcolor{blue}{[F3]}.
\item The kernel of a T3 term is $K_\alpha$ as printed.
\end{itemize}
@@CORRECTED@@

%==========================================================================================
\section{The cancellation, network by network}\label{sec:networks}
Every symmetrized term is written in common variables and simplified with $U^{bc}=U^{cb}$, $U(p)U(p)=1$ and $f^{acd}f^{bcd}=N_c\delta^{ab}$. Within a kernel structure, the colour structures that are equal up to sign form a network; the coefficients of each network are added exactly. Each table lists, for every network:
\begin{itemize}
\item the contribution of every row (in units of $i\NN$ times the network);
\item the sum as written (T1--T3 fixed; red rows included, $\mathcal D$ excluded);
\item the sum after F3 and F4 (red rows excluded, $\mathcal D$ included).
\end{itemize}
\textcolor{red}{I.II$_{\rm F4}$} and \textcolor{red}{I.III$_{\rm F3}$} are the removed terms, and \textcolor{blue}{$\mathcal D$} the added ones. As written, 20 of the 48 networks do not cancel. After the fixes all 40 remaining networks cancel.
@@NETWORKS@@

%==========================================================================================
\section{The operator order}\label{sec:order}
\subsection{Why the order matters for the two-\texorpdfstring{$\rho$}{rho} terms}
\begin{itemize}
\item \emph{c.c.} The cross section is the expectation value of a Hermitian operator. The ``+c.c.'' of the notes is the Hermitian conjugate. With $\rho^\dagger=\rho$ and real $U$, $f$,
\begin{equation}
\big(c\,ABC\big)^\dagger=c^*\,CBA,\qquad CBA=S(ABC)-\tfrac14\big(\{[A,B],C\}+\{[A,C],B\}+\{[B,C],A\}\big)+\dots
\end{equation}
So the c.c.\ has the same symmetrized part as $c^*ABC$, but two-$\rho$ terms of the opposite sign. In your list every c.c.\ term keeps the order of the original term.
\item \emph{Anticommutators.} Group I Row V, Group II Row I and Group III Row III (part 1) contain $\rho\{\rho,\rho\}$ or $\{\rho,\rho\}\rho$ in the notes. Your list writes each as twice one ordering, with twice the coefficient. Then
\begin{equation}
\{B,C\}A-2\,BCA=-[B,C]A=-\tfrac12\{[B,C],A\}+(\text{one }\rho).
\end{equation}
The symmetrized parts agree, but the two-$\rho$ terms differ.
\item \emph{Group I Row II} is not affected by the second point: $f^{abc}\big(\rho^b(x)\rho^c(y)-\rho^b(y)\rho^c(x)\big)=f^{abc}\{\rho^b(x),\rho^c(y)\}$ exactly, so your two orderings are the anticommutator.
\end{itemize}

\paragraph{The rules, row by row} (common variables):
\begin{center}\small
\begin{tabular}{@{}l l l@{}}
row & amplitude ($\alpha,\beta,\gamma$) & c.c.\ ($\alpha',\beta',\gamma'$)\\ \hline
four-$\rho$ remainders A, B & as written ($ABC$) & -- (no c.c.)\\
Group I Row II & $\rho(x')\{\rho(x),\rho(y)\}$ (as written) & $\{\rho(x'),\rho(y)\}\rho(x)$ (reversed)\\
Group I Row III, $\mathcal D$ & as written ($ABC$) & reversed ($CBA$)\\
Group I Row IV & as written & reversed\\
Group I Row V & $A\{B,C\}$, half the coefficient & $\{B,C\}A$, half the coefficient\\
Group II Row I & $\{B,C\}A$ ($\alpha,\beta$), $A\{B,C\}$ ($\alpha',\beta'$), half the coefficient & -- (the two are conjugates)\\
Group III Rows I, II, V, VI & as written & reversed\\
Group III Row III & $A\{B,C\}$, half the coefficient & $\{B,C\}A$, half the coefficient\\
\end{tabular}
\end{center}

\subsection{Literal order versus operator order}
Table~\ref{tab:literal} compares the two-$\rho$ terms of each source obtained by reordering your list literally with those of the operator order. Each value is pointwise at one random configuration.
\begin{table}[H]
\centering\small
\begin{tabular}{@{}l r r l@{}}
source & literal order & operator order & \\ \hline
@@LITERAL@@
\end{tabular}\qquad
\begin{tabular}{@{}r r r r@{}}
seed & total, literal & total, operator & JIMWLK$\otimes$LO\\ \hline
@@LITTOT@@
\end{tabular}
\caption{Left: two-$\rho$ content of each source of three-$\rho$ terms (after the fixes). Right: total two-$\rho$ terms (these sources plus the genuine two-$\rho$ terms and the lists of \texttt{4rho.tex}) against JIMWLK$\otimes$LO. Only the operator order gives JIMWLK$\otimes$LO.}
\label{tab:literal}
\end{table}

\subsection{The corrected list in operator order}\label{sec:oplist}
Your terms after the fixes, with your colour indices, in operator order. Terms whose order or form changed are blue. For an anticommutator the coefficient is half of yours: your single ordering counted it twice. Group I Row II keeps your two orderings per pair, which are exactly $\rho\{\rho,\rho\}$ by the $f^{abc}$ antisymmetry; in its c.c.\ (part 2) each line is reversed.
@@OPLIST@@

%==========================================================================================
\section{The two-\texorpdfstring{$\rho$}{rho} terms}\label{sec:two}
Every term of Sec.~\ref{sec:oplist} is reordered with Eqs.~\eqref{eq:plain}--\eqref{eq:antiA}. The colour algebra is done, and equal terms within each source are collected. The result is a list of entries.
\begin{description}
\item[S:] the kernel structure $S_n$ (Table~\ref{tab:struct});
\item[coeff.:] the coefficient in units of $\NN$;
\item[colour structure:] the colour structure, with the two charges as a symmetric (commuting) product.
\end{description}
Each entry is $\NN\int e^{-ik\cdot(w'-w)}\int_{x,x',z}(\text{coeff.})\times(\text{kernel})\times(\text{colour structure})$. The c.c.\ is included. These are the same as the corresponding entries of \texttt{TwoRho\_all.pdf}, Sec.~5.
\begin{table}[H]
\centering\small
\begin{tabular}{@{}l l@{}}
@@STRUCT@@
\end{tabular}
\caption{The kernel structures of the two-$\rho$ terms.}
\label{tab:struct}
\end{table}

\subsection{From Group I, Row II}
@@TWO:3:I.II@@
\subsection{From Group I, Row III (part 1)}
@@TWO:3:I.III@@
\subsection{From \texorpdfstring{$\mathcal D$}{D}}
@@TWO:D@@
\subsection{From Group I, Row IV}
@@TWO:3:I.IV@@
\subsection{From Group I, Row V}
@@TWO:3:I.V@@
\subsection{From Group II, Row I}
@@TWO:3:II.I@@
\subsection{From Group III, Row I}
@@TWO:3:III.I@@
\subsection{From Group III, Row II}
@@TWO:3:III.II@@
\subsection{From Group III, Row III (part 1)}
@@TWO:3:III.III@@
\subsection{From Group III, Row V}
@@TWO:3:III.V@@
\subsection{From Group III, Row VI}
@@TWO:3:III.VI@@
\subsection{From the four-\texorpdfstring{$\rho$}{rho} remainders \texorpdfstring{$d^3N_a$, $d^3N_b$}{a, b}}
@@TWO:A,B@@

%==========================================================================================
\section{The sum: JIMWLK}\label{sec:sum}
The two-$\rho$ terms of Sec.~\ref{sec:two} are added to two other sets: the genuine two-$\rho$ terms of \texttt{Evolution\_two\_rho.tex} (with fixes F5 and F6), and the two-$\rho$ lists of \texttt{4rho.tex} (with fixes F1 and F2). All of these are written out in \texttt{TwoRho\_all.pdf}. Within each kernel structure, the colour structures are reduced to an independent basis, with exact coefficients that are polynomials in $N_c$.
\begin{itemize}
\item The column ``3$\rho$ list'' is the total from Sec.~\ref{sec:two}.
\item ``genuine'' and ``\texttt{4rho.tex}'' are the two other sets.
\item The last column is JIMWLK$\otimes$LO.
\end{itemize}
All @@NB@@ basis structures with a nonzero entry match. So, with the fixes and the operator order, the two-$\rho$ terms coming from your three-$\rho$ list complete the target evolution: the total coefficient of $\log(\vee/\Lambda)$ is JIMWLK$\otimes$LO. Adding the $\log(\vee/k^+)$ terms of each row gives the coefficient of $\log\vee$, which is BFKL$\otimes$LO (\texttt{TwoRho\_all.pdf}, Sec.~9).
@@SUMS@@

%==========================================================================================
\appendix
\section{Scripts}
\begin{itemize}
\item \texttt{parse.py}: reads \texttt{input/three\_rho\_all.tex} (your list, verbatim) into terms.
\item \texttt{fixes.py}: T1--T3, F3, F4. $\mathcal D$ in your notation is in \texttt{input/D\_terms.tex}.
\item \texttt{check3.py}, \texttt{cmp\_entries.py}: row-by-row and term-by-term comparison with the earlier transcription.
\item \texttt{check3b.py}: Table~\ref{tab:states}.
\item \texttt{networks.py}: Sec.~\ref{sec:networks}.
\item \texttt{oporder.py}: the operator order of every term; check that it reproduces the two-$\rho$ terms of the original expressions.
\item \texttt{two\_cmp.py}: Table~\ref{tab:literal}.
\item \texttt{make\_note.py}: writes this file from \texttt{note\_static.tex}.
\end{itemize}
\end{document}
